\chapter{Offentlig forvaltning og god sagsbehandling}

\chapterlead{Forvaltningsret er den del af retten, der omsætter offentlig magt til konkrete afgørelser. Den beskytter borgeren gennem krav til kompetence, oplysning, upartiskhed, høring, begrundelse og klage.}

\section{Fra lov til afgørelse}

En lov er ofte generel. Den siger eksempelvis, at personer, der opfylder bestemte betingelser, kan få en ydelse, eller at en virksomhed skal have en tilladelse. Forvaltningen skal anvende reglen på en konkret person, virksomhed eller situation. Det kræver både faktumindsamling og juridisk vurdering.

Forvaltningsloven gælder som hovedregel for behandlingen af sager, hvor en forvaltningsmyndighed træffer eller skal træffe en afgørelse, og indeholder også regler om blandt andet inhabilitet, partsaktindsigt, partshøring, begrundelse og klagevejledning \citep{forvaltningsloven}. Offentlighedsloven regulerer en bredere adgang til aktindsigt i den offentlige forvaltning \citep{offentlighedsloven}.

\section{Kompetence}

Før man spørger, om en afgørelse er rimelig, må man spørge, om den rette myndighed har truffet den. Kompetence kan være:
\begin{itemize}
  \item \textbf{Saglig:} Hvilken type spørgsmål må myndigheden tage stilling til?
  \item \textbf{Stedlig:} Hvilket geografisk område eller hvilken kommune er relevant?
  \item \textbf{Personel:} Hvem i organisationen må træffe afgørelsen?
  \item \textbf{Tidsmæssig:} Gælder kompetencen på det tidspunkt, hvor beslutningen blev truffet?
\end{itemize}

En myndighed kan godt have kompetence til at behandle en sag, men ikke til at træffe alle tænkelige beslutninger i den. Det er også nødvendigt at undersøge, om afgørelsen er truffet af den rette person eller efter de krav om delegation, som gælder internt og eksternt.

\section{Inhabilitet og upartiskhed}

En sagsbehandler er inhabil, hvis der foreligger forhold, som kan rejse tvivl om vedkommendes upartiskhed. Det kan være en personlig eller økonomisk interesse, et nært forhold til en part eller tidligere medvirken på en måde, der gør den nye behandling problematisk.

Reglen beskytter ikke kun mod faktisk partiskhed. Den beskytter også tilliden til afgørelsen. Derfor er spørgsmålet ikke altid, om personen faktisk lod sig påvirke, men om situationen udefra giver en rimelig grund til tvivl.

\section{Oplysningsgrundlag og officialprincippet}

Den offentlige myndighed har som udgangspunkt ansvar for, at sagen er tilstrækkeligt oplyst, før den træffer afgørelse. Dette kaldes ofte officialprincippet. Omfanget afhænger af sagstype, indgrebets styrke, oplysningernes tilgængelighed og parternes mulighed for selv at fremskaffe dem.

Et tilstrækkeligt oplysningsgrundlag betyder ikke, at alle tænkelige oplysninger skal indsamles. Det betyder, at myndigheden må kende de fakta, der er nødvendige for at anvende den relevante regel forsvarligt. Et alvorligt indgreb kræver typisk større omhu end en rutinemæssig servicebeslutning.

\begin{examplebox}
\textbf{Eksempel: støtte og indtægt.} Hvis en myndighed skal afgøre, om en borger opfylder et indtægtskrav, må den undersøge den relevante periode, den relevante indkomsttype og eventuelle undtagelser. Det er ikke nok at se på ét tilfældigt kontoudtog, hvis reglen kræver en samlet opgørelse. Omvendt er myndigheden heller ikke forpligtet til at undersøge forhold, der åbenlyst ikke kan påvirke resultatet.
\end{examplebox}

\section{Partshøring}

Partshøring giver parten mulighed for at kommentere væsentlige faktiske oplysninger, som myndigheden vil bruge til ugunst for parten, og som parten ikke kan antages at kende, med de undtagelser der følger af loven. Høringen skal være reel: Parten skal have en rimelig mulighed for at forstå oplysningerne og svare, før afgørelsen træffes.

Høringspligten er ikke det samme som en ret til at bestemme resultatet. Den er en ret til at blive hørt og til at kunne korrigere fejl, forklare kontekst eller fremlægge nye oplysninger. Manglende høring kan få betydning for afgørelsens gyldighed, men virkningen afhænger af den konkrete sag og af, om fejlen kan have haft betydning.

\section{Begrundelse og klagevejledning}

En begrundelse gør afgørelsen forståelig og kontrollerbar. Den bør vise:
\begin{enumerate}
  \item hvilket retsgrundlag myndigheden har brugt,
  \item hvilke hovedhensyn der har været bestemmende ved et skøn, og
  \item hvilke faktiske oplysninger myndigheden har lagt til grund.
\end{enumerate}

En begrundelse er mere end en konklusion. ``Du opfylder ikke betingelserne'' giver sjældent parten mulighed for at vurdere, om myndigheden har forstået sagen rigtigt. En kort afgørelse kan være lovlig, hvis sagen er enkel, men enkelhed skal ikke bruges som undskyldning for at skjule det afgørende ræsonnement.

Når en afgørelse kan påklages, skal klagevejledningen normalt fortælle, til hvem og inden for hvilken frist. Ikke alle afgørelser har en administrativ klageadgang, og den konkrete regel skal derfor kontrolleres. En uklar eller manglende vejledning kan påvirke partens muligheder for at reagere, men den konkrete virkning afhænger af reglerne og omstændighederne.

\section{Skøn, lighed og saglighed}

En regel kan give myndigheden et skøn. Skøn betyder ikke, at myndigheden må vælge helt frit. Skønnet skal holdes inden for lovens formål, bygge på relevante hensyn og respektere lighed og proportionalitet.

\begin{description}[style=nextline,leftmargin=0pt,labelindent=0pt]
  \item[Saglighed] Hensynene skal have forbindelse til den opgave, loven regulerer. Personlige sympatier eller irrelevante politiske ønsker må ikke styre afgørelsen.
  \item[Lighed] Væsentligt ens sager skal behandles ens, medmindre der er en relevant forskel.
  \item[Proportionalitet] Indgrebet skal være egnet, nødvendigt og rimeligt i forhold til formålet.
  \item[Magtfordrejning og skønsgrænser] En myndighed må ikke lade irrelevante hensyn styre eller gøre en hovedregel til en absolut regel, hvis loven forudsætter individuel vurdering.
  \item[Berettigede forventninger] En myndigheds tidligere udmeldinger og praksis kan i visse tilfælde begrænse pludselige ændringer, men forventninger kan ikke uden videre fastholde en ulovlig praksis.
\end{description}

\section{Aktindsigt og tavshed}

Aktindsigt gør det muligt at se, hvad en myndighed ved og har lagt vægt på. Retten er ikke ubegrænset. Hensyn til private forhold, sikkerhed, interne arbejdsdokumenter og andre beskyttede interesser kan begrunde undtagelser. Undtagelser skal fortolkes i lyset af lovens formål, og der kan være pligt til meroffentlighed.

Tavshedspligt beskytter oplysninger, som myndigheden ikke må videregive. Aktindsigt og tavshedspligt skal derfor afvejes i hver enkelt sag. At en oplysning findes hos en myndighed, betyder ikke automatisk, at alle har adgang til den.

\section{Kontrol med forvaltningen}

En afgørelse kan kontrolleres af en intern genvurdering, en administrativ klageinstans, Folketingets Ombudsmand, domstolene eller et særligt nævn. Kontrollen kan være retlig, faglig eller begge dele. En klageinstans kan nogle gange prøve hele sagen, mens domstolen typisk fokuserer på retlige spørgsmål og grænserne for forvaltningens skøn.

\practice{Når du vurderer en myndighedsafgørelse, skriv først en tidslinje. Marker derefter: kompetence, faktagrundlag, høring, hjemmel, skøn, begrundelse, klagevejledning og tidsfrister. Først derefter vurderer du, om resultatet virker rimeligt.}

\question{En myndighed ændrer en borgers ydelse med virkning fra næste måned. Afgørelsen nævner kun et nyt internt notat og skriver, at ``reglerne er blevet strammet''. Hvilke oplysninger og juridiske spørgsmål mangler du?}

\section*{Kapitelopsamling}
\begin{itemize}
  \item En myndighed skal have kompetence og et tilstrækkeligt hjemmelsgrundlag.
  \item Sagen skal være tilstrækkeligt oplyst, og parten skal høres, når reglerne kræver det.
  \item Begrundelsen skal gøre retsgrundlag, faktum og afgørende hensyn kontrollerbare.
  \item Skøn er bundet af saglighed, lighed, proportionalitet og lovens formål.
  \item Aktindsigt, klage og domstolskontrol er centrale retssikkerhedsmekanismer.
\end{itemize}
